\section{Report di Analisi e Revisione: CyTrONE vs
BlueAgent}\label{report-di-analisi-e-revisione-cytrone-vs-blueagent}

\textbf{Documento analizzato:}
\passthrough{\lstinline!1.3.3\_CR\_Alternativi\_CyTrONE.md!}\\
\textbf{Data:} 09/10/2026

\begin{center}\rule{0.5\linewidth}{0.5pt}\end{center}

\subsection{1. Come "blastare" CyTrONE in ottica BlueAgent (Dominio A/D
e Patch
Management)}\label{1-come-blastare-cytrone-in-ottica-blueagent-dominio-ad-e-patch-management}

Il framework \textbf{BlueAgent} è progettato specificamente per
competizioni \textbf{Attack-Defense (A/D)} e scenari complessi di
\textbf{Patch Management}, contesti caratterizzati da vincoli operativi
e architetturali stringenti che CyTrONE non è strutturalmente in grado
di soddisfare.

\subsubsection{Punti critici da evidenziare contro
CyTrONE:}\label{punti-critici-da-evidenziare-contro-cytrone}

\begin{enumerate}
\def\labelenumi{\arabic{enumi}.}
\item
  \textbf{Side-Channel Leakage sistematico nei controlli:}

  \begin{itemize}
  \tightlist
  \item
    In una competizione A/D, la verifica periodica tramite checker
    automatizzati richiede l\textquotesingle invio di exploit/payload
    per testare lo stato di vulnerabilità del servizio.
  \item
    CyTrONE adotta un \emph{cyberattack controller} remoto che scaglia
    attacchi (es. via Metasploit) direttamente attraverso la rete del
    partecipante.
  \item
    \textbf{Conseguenza critica:} Un concorrente può eseguire
    \passthrough{\lstinline!tcpdump!} o Wireshark sulla propria
    interfaccia di rete, intercettare il payload malevolo inviato dal
    checker e \textbf{replicarlo immediatamente contro gli altri team},
    distruggendo l\textquotesingle integrità e l\textquotesingle equità
    della competizione.
  \item
    \textbf{Soluzione BlueAgent:} Separazione netta tra ambiente Live e
    Sandbox. Le verifiche avvengono su cloni Copy-on-Write (CoW) isolati
    tramite OVN/Linux network namespaces, impedendo qualunque sniffing o
    fuga di payload.
  \end{itemize}
\item
  \textbf{Rischio di Denial of Service / Crash del sistema Live:}

  \begin{itemize}
  \tightlist
  \item
    L\textquotesingle esecuzione di payload aggressivi o malformati
    contro l\textquotesingle istanza live di training può provocare
    kernel panic, corruzione di database o blocco dei servizi.
  \item
    In CyTrONE ciò compromette l\textquotesingle intero ambiente del
    partecipante, richiedendo interventi manuali o lunghi reset.
  \item
    In BlueAgent, il test viene eseguito esclusivamente su
    un\textquotesingle istanza clonata effimera: anche in caso di crash
    totale del clone, l\textquotesingle ambiente di produzione del team
    rimane intatto e operativo.
  \end{itemize}
\item
  \textbf{Latenza di Provisioning e Ripristino (Rollback
  Inaccettabile):}

  \begin{itemize}
  \tightlist
  \item
    Durante una CTF è frequente che un concorrente renda instabile il
    proprio sistema e necessiti di un ripristino.
  \item
    CyTrONE impiega minuti (fino a 8+ minuti per scenari base con sole 2
    VM) e degrada esponenzialmente con l\textquotesingle aumentare dei
    nodi.
  \item
    BlueAgent garantisce provisioning e rollback in
    \textbf{\textasciitilde100--500 ms} grazie all\textquotesingle uso
    di container di sistema Incus/LXC e snapshot a livello di blocco
    (ZFS/Btrfs).
  \end{itemize}
\item
  \textbf{Inadeguatezza dell\textquotesingle Isolamento Topologico:}

  \begin{itemize}
  \tightlist
  \item
    Una CTF A/D moderna richiede topologie di rete complesse e omogenee
    (es. subnet identiche \passthrough{\lstinline!10.k.1.0/24!} per ogni
    team \(k\)) per non dover riscrivere script o regole di routing.
  \item
    CyTrONE si affida a clonazione tradizionale
    (\passthrough{\lstinline!parallel-scp!}) su hypervisor KVM,
    risultando rigido e non scalabile nella gestione di reti
    virtualizzate multi-tenant isolate.
  \end{itemize}
\end{enumerate}

\begin{center}\rule{0.5\linewidth}{0.5pt}\end{center}

\subsection{2. Punti da Approfondire nel
Testo}\label{2-punti-da-approfondire-nel-testo}

\begin{itemize}
\tightlist
\item
  \textbf{Quantificazione dell\textquotesingle overhead delle Macchine
  Virtuali (VM):}

  \begin{itemize}
  \tightlist
  \item
    Esplicitare i costi reali: 1--4 GB di RAM pre-allocata per VM guest,
    overhead di virtualizzazione CPU (VT-x/AMD-V) e boot time completo
    del kernel OS.
  \item
    Confrontarlo con i container di sistema (LXC/Incus): overhead di
    memoria ridotto a \textasciitilde100--300 MB e condivisione sicura
    del kernel host.
  \end{itemize}
\item
  \textbf{Dicotomia "Training Guidato" vs "Competizione Dinamica":}

  \begin{itemize}
  \tightlist
  \item
    CyTrONE assume verifiche elementari basate sul confronto di
    stringhe/flag note (\emph{Assunzioni 1 e 2}).
  \item
    Nel Patch Management contemporaneo la validazione è \textbf{stateful
    e bivalente}: deve attestare sia la neutralizzazione della falla,
    sia la conservazione dell\textquotesingle integrità funzionale (SLA
    del servizio per utenti legittimi).
  \end{itemize}
\end{itemize}

\begin{center}\rule{0.5\linewidth}{0.5pt}\end{center}

\subsection{3. Elementi Mancanti di Rilievo
Architetturale}\label{3-elementi-mancanti-di-rilievo-architetturale}

\begin{itemize}
\tightlist
\item
  \textbf{Mancanza del concetto di Sandbox / Dual-Environment:}

  \begin{itemize}
  \tightlist
  \item
    CyTrONE opera su una singola istanza per utente. Manca completamente
    il paradigma fondamentale di BlueAgent: \textbf{biforcazione tra
    Live Environment e Sandbox effimera}.
  \end{itemize}
\item
  \textbf{Modernità dello Storage (CoW vs Immagini monolitiche):}

  \begin{itemize}
  \tightlist
  \item
    La clonazione di intere immagini disco via rete con
    \passthrough{\lstinline!parallel-scp!} è un collo di bottiglia non
    accettabile; mancano riferimenti alla gestione moderna tramite
    filesystem Copy-on-Write (ZFS/Btrfs).
  \end{itemize}
\item
  \textbf{Sicurezza di Connessione e Anti-Spoofing:}

  \begin{itemize}
  \tightlist
  \item
    Assenza di meccanismi di VPN crittograficamente confinata (come il
    Cryptokey Routing di WireGuard in BlueAgent) per prevenire
    intercettazioni e spoofing tra team differenti.
  \end{itemize}
\end{itemize}

\begin{center}\rule{0.5\linewidth}{0.5pt}\end{center}

\subsection{4. Revisione Qualitativa del Documento e Correzioni
Formali}\label{4-revisione-qualitativa-del-documento-e-correzioni-formali}

\subsubsection{Punti di Forza:}\label{punti-di-forza}

\begin{itemize}
\tightlist
\item
  Ottima linearità nell\textquotesingle esposizione dei macro-componenti
  (CyLMS, CyRIS, CyPROM).
\item
  Sezione benchmark quantitativi (1.3.3.1) chiara, efficace e ben
  argomentata tramite dati sperimentali (StarBED).
\item
  Perfetta impostazione per introdurre il contrasto prestazionale
  cercato dalla tesi.
\end{itemize}

\subsubsection{Correzioni Formali da apportare (Typo e
Sintassi):}\label{correzioni-formali-da-apportare-typo-e-sintassi}

\begin{enumerate}
\def\labelenumi{\arabic{enumi}.}
\tightlist
\item
  \textbf{Riga 7:} \passthrough{\lstinline!si occupa d icaricare!}
  \(\rightarrow\) \passthrough{\lstinline!si occupa di caricare!}
\item
  \textbf{Riga 14:} \passthrough{\lstinline!si conentra!}
  \(\rightarrow\) \passthrough{\lstinline!si concentra!}
\item
  \textbf{Riga 16:}
  \passthrough{\lstinline!Quindi l'esperienza utente si concentra sul fatto che...!}
  \(\rightarrow\) Stile informale. Riformulare in:\\
  \emph{«L\textquotesingle esperienza utente è incentrata sulla
  fruizione del percorso didattico tramite piattaforma LMS e sul
  contestuale accesso al Cyber Range dedicato.»}
\item
  \textbf{Riga 23:} \passthrough{\lstinline!Training Coordinator".Le!}
  \(\rightarrow\) Spaziatura errata dopo le virgolette e il punto
  (\passthrough{\lstinline!Training Coordinator". Le!}).
\item
  \textbf{Riga 31:} \passthrough{\lstinline!vulnearbilita!}
  \(\rightarrow\) \passthrough{\lstinline!vulnerabilità!}
\item
  \textbf{Riga 33 e 35:} \passthrough{\lstinline!Tutta via!}
  \(\rightarrow\) \passthrough{\lstinline!Tuttavia!}
\item
  \textbf{Riga 59:}
  \passthrough{\lstinline!Nella figura sborgz viene analizzato...!}
  \(\rightarrow\) \textbf{Placeholder rimasto nel testo
  (\passthrough{\lstinline!sborgz!})}. Sostituire con il corretto
  riferimento (\passthrough{\lstinline!Nella Figura X!} o
  \passthrough{\lstinline!Nella figura sottostante!}).
\item
  \textbf{Riga 66:}
  \passthrough{\lstinline!Con il nostro lavoro di tesi, si vuole cercare di abbattere...!}
  \(\rightarrow\) Preferibile formula accademica attiva:\\
  \emph{«L\textquotesingle obiettivo del presente lavoro di tesi è
  abbattere il tempo complessivo di approvvigionamento \(T\),
  riducendolo dall\textquotesingle ordine dei minuti a quello dei
  secondi su hardware non dedicato.»}
\end{enumerate}
